\questao{111}{corD}{Energia e calor}{Fontes de energia: renováveis e não renováveis}{Só de conceito}
\begin{enunciado}
% gerado por ferramentas/md2tex_enunciado.py a partir de questoes/q111.md (não editar à mão)
A busca por fontes alternativas de energia é uma das grandes preocupações da humanidade devido ao possível esgotamento dos combustíveis fósseis e à degradação ambiental causada por eles. Como a demanda de energia cresce significativamente, há um grande investimento em outros tipos de geração de energia, tais como hidrelétrica, nuclear, solar, termelétrica a biomassa e eólica.\par
\fonte{BIZERRA, Ayla M. C.; QUEIROZ, Jorge L. A.; COUTINHO, Demétrios A. M. O impacto ambiental dos combustíveis fósseis e dos biocombustíveis: as concepções de estudantes do ensino médio sobre o tema. \textbf{Revbea}, n. 3, 2018 (adaptado).}
\figq[0.78\linewidth]{q111_fig1.png}
Dentre as fontes de energia citadas, o tipo de energia oriunda de uma fonte não renovável, que independe das condições climáticas para ser obtida, é a\par
\begin{alternativas}
  \item nuclear, gerada pela fissão nuclear do urânio.
  \item hidrelétrica, gerada pela movimentação da água.
  \item solar, gerada pelo uso de células fotovoltaicas.
  \item termelétrica, produzida pela combustão de biomassa.
  \item eólica, produzida pelo movimento das pás de aerogeradores.
\end{alternativas}

\end{enunciado}

\quer{qual das cinco fontes de energia é, ao mesmo tempo, \textbf{não renovável} e \textbf{independente do clima}.}

\begin{ideia}
A questão impõe \textbf{duas condições}. A resposta precisa cumprir as duas ao mesmo tempo.
\begin{itemize}
  \item \textbf{Fonte renovável:} a natureza repõe em pouco tempo. Sol, vento, chuva e plantas (que podem ser replantadas).
  \item \textbf{Fonte não renovável:} existe em quantidade limitada e um dia acaba. Petróleo, carvão mineral, gás natural e
        \textbf{urânio} (um minério retirado de minas).
  \item \textbf{Depende do clima:} só funciona bem se o tempo ajudar (sol, vento ou chuva).
\end{itemize}
A melhor estratégia é montar um quadro e testar as duas condições em cada fonte.
\end{ideia}

\passoapasso
\npasso{Teste as duas condições em cada fonte.}
{\small\renewcommand{\arraystretch}{1.35}
\begin{center}
\begin{tabular}{@{}>{\raggedright\arraybackslash}p{3.2cm}>{\raggedright\arraybackslash}p{5.4cm}>{\raggedright\arraybackslash}p{6.3cm}c@{}}
\toprule
\textbf{Fonte} & \textbf{É não renovável?} & \textbf{Independe do clima?} & \textbf{Serve?}\\
\midrule
\alt{A} Nuclear & \textbf{Sim.} O urânio é um minério que se esgota. & \textbf{Sim.} A usina funciona de dia, de noite, com chuva ou com seca. & \textbf{Sim}\\
\alt{B} Hidrelétrica & Não. A chuva repõe a água dos rios. & Não. Na seca, os reservatórios baixam. & Não\\
\alt{C} Solar & Não. O Sol brilha todos os dias. & Não. Nuvens e a noite reduzem a geração. & Não\\
\alt{D} Termelétrica a biomassa & Não. As plantas (cana, madeira) podem ser replantadas. & Em boa parte, sim: a biomassa pode ser guardada e queimada quando for preciso. & Não\\
\alt{E} Eólica & Não. O vento não acaba. & Não. Sem vento, as pás param. & Não\\
\bottomrule
\end{tabular}
\end{center}}

\npasso{Escolha a única que cumpre as duas condições.}
Só a energia \textbf{nuclear} tem ``sim'' nas duas colunas.

\resposta{A}{A energia nuclear usa urânio, que é um recurso limitado (não renovável), e a usina funciona com qualquer tempo.}

\begin{ideia}[Como uma usina nuclear gera eletricidade]
Na \textbf{fissão nuclear}, o núcleo de um átomo de urânio se parte em dois núcleos menores e libera muita energia em forma de calor
(não confunda com a \emph{fusão} da questão 94, em que dois núcleos se juntam).
Esse calor ferve água. O vapor gira uma turbina. A turbina gira um gerador, que produz a energia elétrica.
Do vapor em diante, é o mesmo esquema de uma termelétrica comum: o que muda é a origem do calor.
\end{ideia}

\begin{pegadinha}[Onde muita gente erra]
\begin{itemize}
  \item \textbf{Marcar D por causa da palavra ``termelétrica''.} Muita gente pensa logo em carvão e petróleo. Mas o enunciado diz
        termelétrica \textbf{a biomassa}, que é renovável.
  \item \textbf{Olhar só uma condição.} A hidrelétrica, a solar e a eólica falham nas duas. A biomassa falha em uma.
        Só vale a fonte que cumpre as duas.
  \item \textbf{Confundir ``limpa'' com ``renovável''.} A energia nuclear não solta fumaça, mas o seu combustível é limitado.
\end{itemize}
\end{pegadinha}
